\section{From a structural model to an identified simulator}\label{from-a-structural-model-to-an-identified-simulator}

\subsection{1. Structural parameter blocks}\label{structural-parameter-blocks}

The paper partitions parameters as

\[
\theta=(\theta_G,\theta_B,\theta_M,\theta_C),
\]

where:

\begin{itemize}
\tightlist
\item
  \texttt{theta\_G} controls growth, mortality, and quality transition;
\item
  \texttt{theta\_B} controls harvest, taper, bucking, and tree-to-log recovery;
\item
  \texttt{theta\_M} controls processor acceptance, matching, price response, and
  bargaining;
\item
  \texttt{theta\_C} controls service, verification, and operating cost.
\end{itemize}

This decomposition is useful because each block requires different evidence.
A model can be strongly identified in one block and weakly identified in
another.

\subsection{2. Observation model with discrepancy}\label{observation-model-with-discrepancy}

The paper writes the generic observation equation as

\[
Y=F_{obs}(X,u,\theta)+\delta_{model}+\epsilon.
\]

Here:

\begin{itemize}
\tightlist
\item
  \texttt{F\_obs} is the structural observation map;
\item
  \texttt{delta\_model} is model discrepancy;
\item
  \texttt{epsilon} is measurement noise.
\end{itemize}

The distinction matters. Suppose measured DBH growth differs systematically
from model predictions at hot, dry sites. If discrepancy is omitted, a Bayesian
fit may force growth parameters into values that compensate for a missing
mechanism. The posterior can then become precise around the wrong structural
story.

Model discrepancy requires regularization and experimental designs that separate
parameter effects from structural error. These restrictions preserve
falsifiability.

This treatment follows the calibration framework of \citet{kennedy2001calibration}.

\subsection{3. Data-to-parameter identification map}\label{data-to-parameter-identification-map}

The working paper proposes the following hierarchy.

\subsubsection{Tree and plot remeasurement}\label{tree-and-plot-remeasurement}

Repeated DBH, height, survival, and quality/form observations inform
$\theta_G$. Repeated individual-tree measurements are especially useful because
increments can be conditioned on prior state. With only two time points, process
diffusion, measurement noise, and latent heterogeneity can be hard to separate.

\subsubsection{Harvest and bucking records}\label{harvest-and-bucking-records}

Tree-level pre-harvest measurements linked to recovered log dimensions and grade
outcomes inform $\theta_B$. In their absence, the mature-tree log-recovery kernel
remains uncertain because standing volume supplies insufficient information.

\subsubsection{Contractor records}\label{contractor-records}

Time, equipment, fuel, payload, distance, labour, and realized productivity
inform the delivered-cost kernel and operator state.

\subsubsection{Processor intake and grading}\label{processor-intake-and-grading}

Observed log dimensions, accepted/rejected status, realized grade, moisture,
quality, and price inform processor grade maps and measurement-error-aware
acceptance.

\subsubsection{Repeated bids and alternative buyers}\label{repeated-bids-and-alternative-buyers}

These help separate technical delivered surplus from bargaining and outside
options.

\subsubsection{Product conversion and renewal}\label{product-conversion-and-renewal}

Observed customer adoption, retention, and response to price inform the
reduced-order demand model.

\subsection{4. Identifiability is a property of the experiment}\label{identifiability-is-a-property-of-the-experiment}

Even a mathematically identifiable parameter can be practically unidentified
under a weak data design. If all observed stands have nearly the same density,
for example, a competition coefficient may be confounded with site
productivity. If all logs are far from a grade boundary, measurement-error scale
in the grade classifier may have little support.

The control framework turns this into an experimental-design problem. A useful
experiment is one that creates variation where competing parameter explanations
make different predictions.

\subsection{5. Particle approximation of the biological law}\label{particle-approximation-of-the-biological-law}

For a site/mark class, simulate particles

\[
X_t^{1,N},\ldots,X_t^{N,N}
\]

using the same common-noise path and independent idiosyncratic noise. The
empirical law approximates the conditional distribution.

A practical implementation should preserve:

\begin{itemize}
\tightlist
\item
  shared common shock within a site-year realization;
\item
  individual heterogeneity;
\item
  mortality and harvest as mass-changing events;
\item
  weights if particles represent unequal numbers of real trees.
\end{itemize}

The finite particle system is also the natural place to retain tree labels for
synthetic trajectory diagnostics. The limiting measure remains unlabelled.

\subsection{6. Discretizing log space}\label{discretizing-log-space}

After harvest and bucking, the model produces a measure over end diameters,
length, quality, moisture, biological class, and location.

A uniform grid can waste resolution far from processor thresholds and be too
coarse near them. Adaptive binning is preferable where value is discontinuous or
steep.

One useful hierarchy is:

\begin{Shaded}
\begin{Highlighting}[]
\NormalTok{individual tree particles}
\NormalTok{    {-}\textgreater{} stochastic bucking/recovery}
\NormalTok{    {-}\textgreater{} adaptive log bins}
\NormalTok{    {-}\textgreater{} processor{-}grade nodes}
\NormalTok{    {-}\textgreater{} operator{-}capacity nodes}
\end{Highlighting}
\end{Shaded}

Binning only by tonnes is insufficient if processor value changes sharply with
log dimensions.

\subsection{7. Exact LP as a matching benchmark}\label{exact-lp-as-a-matching-benchmark}

The finite discretization of the allocation problem is a linear program. Solve
this exact LP on representative scenarios to establish a reference value and
reference dual potentials.

For larger problems, entropy-regularized transport can trade accuracy for speed.
The approximation should be diagnosed by decreasing the regularization parameter
and comparing:

\begin{itemize}
\tightlist
\item
  primal objective value;
\item
  matched grade quantities;
\item
  dual potentials;
\item
  sensitive edges near compatibility thresholds.
\end{itemize}

The regularized solver must preserve the scarcity signals used by the outer
control problem.

\subsection{8. Dual values as a sensitivity oracle}\label{dual-values-as-a-sensitivity-oracle}

Suppose an outer decision changes available supply, demand, or capacity
approximately affinely. Because the inner matching value is concave in these
resources, optimal dual variables act as supergradients.

This enables decomposition methods. The outer problem can propose a sensing or
coverage decision; the inner LP returns both value and marginal-value
information; those dual values can generate cuts or local approximations for the
outer search.

\subsection{9. Posterior approximation}\label{posterior-approximation}

The exact joint posterior over physical measures and parameters is intractable
for a large system. Candidate approximations include:

\begin{itemize}
\tightlist
\item
  sequential Monte Carlo over parameter and state particles;
\item
  scenario trees;
\item
  stochastic model-predictive control;
\item
  rollout from a low-dimensional belief embedding;
\item
  fitted value functions on decision-relevant summaries.
\end{itemize}

The reduction should preserve variables that affect decisions. Useful candidates
from the structural model include:

\begin{itemize}
\tightlist
\item
  local compatible supply by grade and radius;
\item
  matched-demand fractions;
\item
  posterior moments of size and quality;
\item
  high-value dual potentials;
\item
  uncertainty near specification boundaries.
\end{itemize}

Raw hectares remain a useful exposure variable. Grade-resolved supply, matched
demand, and posterior uncertainty provide additional belief coordinates.

\subsection{10. A nested simulation path}\label{a-nested-simulation-path}

For posterior draw \texttt{s}:

\begin{enumerate}
\def\labelenumi{\arabic{enumi}.}
\tightlist
\item
  sample structural parameters and a common environmental path;
\item
  propagate tree particles;
\item
  simulate mortality and harvest;
\item
  optimize or sample bucking;
\item
  convert recovered logs to mass using state-dependent density/volume;
\item
  construct feasible processor/operator edges;
\item
  solve matching and retain primal/dual outputs;
\item
  simulate observations generated by the chosen action;
\item
  update the posterior;
\item
  evaluate risk and continuation value.
\end{enumerate}

This procedure defines a stochastic program over posterior and process
distributions.

Scenario-count diagnostics can use sample-average approximation bounds
\citep{kleywegt2002saa}. Correlated input sensitivity can be summarized with
Shapley effects \citep{iooss2019shapley}.

\subsection{11. Decompose numerical uncertainty}\label{decompose-numerical-uncertainty}

At least five error sources should be tracked separately.

\subsubsection{Biological discretization error}\label{biological-discretization-error}

Finite particle number, time step, and state approximation.

\subsubsection{Filter/posterior approximation error}\label{filterposterior-approximation-error}

Finite posterior particles, approximate likelihoods, or belief compression.

\subsubsection{Recovery model error}\label{recovery-model-error}

Taper, defect, bucking, and log-quality uncertainty.

\subsubsection{Matching/scenario error}\label{matchingscenario-error}

Discretization of log space, entropy regularization, and finite scenario count.

\subsubsection{Structural discrepancy}\label{structural-discrepancy}

Real mechanisms missing from the model family.

Additional Monte Carlo samples reduce sampling error and leave structural
discrepancy unchanged.

\subsection{12. Minimum reproducibility standard}\label{minimum-reproducibility-standard}

A numerical experiment should record:

\begin{itemize}
\tightlist
\item
  code revision;
\item
  parameter source and whether synthetic, empirical, or prior;
\item
  random seed where relevant;
\item
  particle count and time step;
\item
  discretization/bins;
\item
  matching solver and tolerances;
\item
  regularization value if used;
\item
  posterior approximation method;
\item
  output uncertainty by source where feasible.
\end{itemize}

This is the bridge from mathematical theory to a research-grade implementation.
